
\chapter{Diagonalisation}

\justify

\setlength{\parindent}{0pt}

% ==================================================================================================================================
% Introduction

Lors de la résolution de gros systèmes linéaires, il serait intéressant de "modifier" ce systèmes pour en trouver un équivalent plus simple, 
ce qui faciliterai sa résolution. 
En écrivant un système linéaire comme une équation matrice/vecteur, on pourrait donc trouver une nouvelle base de représentation de 
la l'endormorphisme associé dans laquelle sa matrice serait diagonale/triangulaire. 

Dans ce chapitre, nous allons voir quelques critères permettant de faire ce changement de base et les méthodes de recherche des dites matrices/bases. 

\begin{remark}[Notation]
    Dans ce chapitre, nous noterons $f$ des endomorphismes et $M$ les matrices représentant ces endormorphismes. 
    Il est important de comprendre la différence entre un endomorphisme sa matrice représentative. 
    En effet, un endormorphisme est une application linéaire et sa matrice représentative permet de le caractériser dans une certaine base. 
    Lorsque l'on change de base (i.e de point de vue), l'endomorphisme reste le même mais sa matrice représentative peut changer. 
\end{remark}

% ==================================================================================================================================
% Element Propres 

\section{Eléments Propres}

\begin{definition}[Vecteur et valeur propre]
    Soit $ f \in \mathcal{L}(E)$, $v \in E$ est un \textbf{vecteur propre} de $f$ associé à la \textbf{valeur propre} $\lambda$ ssi :
    \[ 
        \begin{cases}
            v \not = 0 \\
            f(v) = \lambda v
        \end{cases}
    \]
\end{definition}

Il est important de noter qu'un vecteur propre est toujours associé à une valeur propre. 

\begin{definition}[Sous-espace propre]
    On appelle \textbf{sous-espace propre} de f (représenté par la matrice M) associé à la valeur propre $\lambda$ le sous-espace vectoriel de $E$ : 
    \[ \boxed{E_\lambda = ker(M - \lambda Id_E)} \]
\end{definition}

On remarquera la praticité de l'utilisation du théorème du rang pour déterminer la dimension du sous-espace propre : \[ dim(E_\lambda) = dim(E) - rg(M - \lambda Id_E) \]

\begin{definition}[Polynôme Caractéristique]
    Soit $A \in \mathcal{M}_n(\mathbb{K}) $. On appelle \textbf{polynôme caractéristique} de $A$ le polynôme : 
    \[ P_A = det(A - X \times Id_E) \]
\end{definition}

Le terme "polynôme caractéristique" d'une matrice/d'un endomorphisme prend tout son sens. 
En effet, il permet vraiment de caractériser un endormorphisme par sa matrice représentative dans une base. 

\begin{remark}
    Le polynôme caractéristique d'un endomorphisme $f$ est dit \textbf{intrinsèque}.
    (i.e le polynôme caractéristique de $f$ ne dépend pas du choix de la matrice représentant de $f$)
\end{remark}



% ==================================================================================================================================

\section{Définitions et caractérisations}

On cherche ici à savoir si un endormorphisme ou une matrice peut être représenté par une matrice diagonale dans une certaine base 
et à avoir une méthode pour construire une telle base. 

\begin{definition}[Endomorphisme diagonalisable]
    Un endomorphime $f$ est dit \textbf{diagonalisable} ssi il existe une base $E$ dans laquelle il est représenté par une matrice $D$ diagonale. 
    Une matrice $A$ est dite diagonalisable ssi elle est semblable à une matrice diagonale. 
    \[ \boxed{ \text{A est diagonalisable} \; \Leftrightarrow \; \exists P \in GL_n(\mathbb{K}), \; A = PDP^{-1} } \]
\end{definition}

\begin{proposition}
    Un endomorphisme $f$ est diagonalisable dans E ssi il existe une base $\mathcal{B}$ de E 
    constituée des vecteurs propres de $f$. 
    La matrice diagonale associé sera donc composée des valeurs propres associés. 
\end{proposition}

\begin{proposition}
    Un endomorphime $f$ représenté par une matrice $M$ est diagonalisable ssi M l'est. 
\end{proposition}

% La proposition ci-dessus permet donc de faire le lien entre une matrice et un endomorphisme diagonalisable. 

\begin{proposition}
    Soit f un endomorphisme représenté par une matrice A dans une base $\mathcal{B}$ de dimension n. 
    Soit $ \lambda \in \R^*$. 
    \begin{align*}
        \lambda \text{ est valeur propre de } f & \iff E_\lambda \not = \{ 0 \} \\ 
                                                & \iff f - \lambda \text{Id} \text{ n'est pas de rang } n \\ 
                                                & \iff A - \lambda \text{Id} \text{ n'est pas de rang } n \\ 
                                                & \iff A - \lambda \text{Id} \text{ n'est pas inversible } \\ 
                                                & \iff \det (A - \lambda \text{Id}) = 0 
    \end{align*}
\end{proposition}

Cette proposition nous donne donc une caractérisation très pratique des valeurs propres. 
En plus du théorème du rang, on dispose donc déjà d'outils puissant pour étudier la diagonalisation d'une matrice/d'un endomorphime. 


\section{Recherche des éléments propres}

En pratique, pour diagonaliser un endormorphisme, il nous faut donc commencer par trouver ses éléments propres. 
D'après la proposition précédente, une fois les valeurs propres trouvées, il reste à déterminer les sous-espaces propres $E_\lambda$.
On se ramène ici à calculer des noyaux par le pivot de Gauss. 

\begin{prop}[Lien Valeurs Propres/Racines]
    Soient $A \in \mathcal{M}_n(\K)$ et $\lambda \in \K$ une valeur propre de A. 
    Alors $\lambda$ est une racine de $P_A$. 
    Autrement dit, les valeurs propres d'une matrice sont exactement les racines de son polynôme caractéristique. 
\end{prop}

\begin{definition}[Multiplicité]
    Soit $f$ un endormorphisme représenté par une matrice $A$ dans un base $\mathcal{B}$. 
    La multiplicité d'une valeur propre $\lambda \in \K$ de $f$ est exactement la multiplicité de $\lambda$ en tant que racine de $P_A$.     
\end{definition}

\begin{example}[Recherche des éléments propres]
    Soit $A \in \mathcal{M}_3(\R)$ telle que 
    \[A = 
        \begin{pmatrix}
            2 & 1 & 1 \\ 
            1 & 2 & 1 \\ 
            0 & 0 & 3
        \end{pmatrix}
    \]
    Déterminons les éléments propres de A. Commençons par calculer son polynôme caractéristique. 

    \begin{align*}
        P_A &=  \begin{vmatrix}
                    2-X & 1 & 1 \\ 
                    1 & 2-X & 1 \\ 
                    0 & 0 & 3-X 
                \end{vmatrix}
            = (3-X)    \begin{vmatrix}
                            2-X & 1  \\ 
                            1 & 2-X
                        \end{vmatrix} \\ 
            &= (3-X)((2-X)^2-1) \\
            &= (3-X)(4 - 4X + X^2 - 1) \\
            &= (3-X)^2(1-X)
    \end{align*}

    D'où : $Sp(A) = \{3,1\}$ et 3 est de multiplicité 2. 
\end{example}


\section{Critères de Diagonalisation}

\subsection{Stabilité et somme directe}

\begin{definition}[Sous espace stable]
    Un s.e.v $F$ de $E$ est dit stable par $f$ ssi :
    \[ \forall v \in F, f(v) \in F \quad \text{i.e} \quad f(F) \subset F \]
\end{definition}

\begin{prop}[Sous-espace propre et stabilité]
    Un sous espace propre $E_\lambda$ de f est stable par f. 
\end{prop}

\begin{definition}[Somme Directe]
    Soit E un espace vectoriel et $(F_i)_{i \in I}$ une famille finie de sous-espaces vectoriels de E. 
    On dit que les $F_i$ sont en somme directe ssi pour tout vecteur $v \in \sum_{i \in I} F_i$, il existe une \textbf{unique} 
    décomposition de $v$ comme somme de vecteurs $v_i \in F_i, \forall i \in I$. 
\end{definition}

\begin{prop}[Somme Directe]
    Toute somme de sous espaces propres d'un même endomorphime est directe. 
\end{prop}


\subsection{Caractérisation d'un endormorphisme diagonalisable}

\begin{proposition} Deux propositions très utiles pour déterminer si un endomorphime est diagonalisable : 
    Un endomorphisme $f$ est diagonalisable ssi ses sous espaces propres sont supplémentaires dans $E$ $i.e$ :
        \[ E = \bigoplus_{\lambda \in Sp(f)} E_\lambda \quad \iff \quad \sum_{\lambda \in Sp(f)} dim(E_\lambda) = n \]    
\end{proposition}

Le second point de la proposition précédente nous donne une approximation de la dimension d'un sous-espace propre. 

\begin{prop}[Multiplicité et encadrement]
    Si $\lambda$ est valeur propre de $f$ de multiplicité $\alpha$ alors : $$ 1 \leq \text{dim} (E_\lambda) \leq \alpha $$ 
\end{prop}

On peut donc déterminer un critère fondamental de diagonalisation d'endomorphismes (matrices). 

\begin{criteria}[Critère Fondamental de Diagonalisation]
    Un endomorphime $f$ est diagonalisable ssi 
    \begin{itemize}
        \item son polynôme caractéristique est scindé sur $\K$ et pour chaque valeur propre, 
        la dimension du sous-espace propre associé est égale à la multiplicité de la valeur propre. 
        \item la somme de la dimension de ses sous-espaces propres est $n$. 
    \end{itemize}
\end{criteria}

\begin{corollary}[Critère de Diagonalisation]
    Si $f$ admet n valeurs propres distinctes alors $f$ est diagonalisable. 
\end{corollary}


\section{Méthode Pratique de Diagonalisation}

Maintenant que nous savons déterminer si un endomorphisme est diagonalisable, mettons en place 
une méthode pratique nous donnant la base de diagonalisation, les sous-espaces propres, et la matrice diagonale. 

Soit $f$ un endomorphisme représenté par une matrice $A$ dans une base $\mathcal{B}$. 
Pour diagonaliser $f$, il suffit de :
\begin{enumerate}
    \item Calculer le polynôme caractéristique de $A$. 
    \item Factoriser le polynôme trouvé pour déterminer le spectre de $f$. 
    \item Pour chaque valeur propre $\lambda$ distincte de $f$ :
        \begin{enumerate}
            \item Calculer $A - \lambda \; Id$. 
            \item Echelonner la matrice obtenue par le pivot de Gauss. Avec mémorisation. 
            \item Extraire du pivot une base du sous-espace. 
        \end{enumerate}
    \item Maintenant on peut se trouver face à deux situations :
        \begin{itemize}
            \item Si pour chaque valeur propre, la dimension du sous-espace engendré est égale à la multiplicité de 
                celle-ci, alors $f$ est digonalisable. On poursuit alors la diagonalisation. 
            \item Sinon, $f$ n'est pas diagonalisable. On s'arrête ici... 
        \end{itemize}
    \item La base de diagonalisation sera donc l'union des bases des sous-espaces propres. 
            La matrice diagonale sera composée des valeurs propres de $f$. 
            Pour effectuer le changement de base, on veillera bien à inverser la matrice de passage. 
\end{enumerate}

\begin{example}
    Ici un exemple pratique. 
\end{example}

\section{Introduction à la trigonalisation}

La diagonalisation semble être un outil très fort pour effectuer toute sorte de calculs sur des matrices/endomorphismes. 
Mais en pratique, on se rend compte que très peut de matrices sont diagonalisables. 
On peut donc trouver une autre forme de réduction d'un endomorphisme représenté par une matrice mais, cette fois ci, pas sous forme triagonale, 
mais triangulaire. On va ainsi définir une nouvelle forme de réduction et trouver des critères pour savoir si une matrice est trigonalisable. 

% \subsection{Définitions}

\begin{definition}[Trigonalisation]
    Un endomorphisme $f$ représenté par une matrice A dans une base $\mathcal{B}$ est dit trigonalisable si 
    il existe une base $\mathcal{B}'$ dans laquelle il est représenté par une matrice triangulaire supérieure. 

    \vspace{0.3cm}

    De même une matrice $A$ est dite trigonalisable si elle est semblable à une matrice triangulaire supérieure. 
\end{definition}

\begin{proposition}
    Un endomorphime $f$ représenté par une matrice $A$ est trigonalisable ssi $A$ l'est. 
\end{proposition}

% \subsection{Premiers critères}

\begin{criteria}[Fondamental de Trigonalisation]
    Soient E un $\K$-espace vectoriel et $f \in \mathcal{L}(E)$. $f$ est trigonalisable ssi son polynôme caractéristique est scindé sur $\K$. 
\end{criteria}

\begin{corollary}[La puissance de $\C$]
    Tout endormorphisme est trigonalisable sur $\C$. 
\end{corollary}

On remarque que l'élaboration d'une méthode exaustive de trigonalisation de matrices semble plus complexe que la diagonalisation. 
Nous en verrons une dans le chapitre dédié, mais pour cela, il nous faut définir de nouveaux object qui vont nous permettre d'aller plus loin dans la théorie de la réduction. 

